\documentclass[11pt]{article}
\usepackage[italian]{babel}
\usepackage[utf8]{inputenc}
\usepackage{amsmath}
\usepackage{amssymb}
\usepackage{amsthm}
\usepackage{xspace}
\usepackage{geometry}
\geometry{a4paper, top=1.9cm, bottom=1.9cm, left=2.2cm, right=2.2cm}

\usepackage[colorlinks,allcolors=black]{hyperref}

% Comandi di comodo
\newcommand{\N}{\mathbb{N}}
\newcommand{\R}{\mathbb{R}}
\newcommand{\Z}{\mathbb{Z}}
\newcommand{\E}{\mathbb{E}}
\newcommand{\val}{\operatorname{val}}
\newcommand{\CSP}{\mathtt{CSP}\xspace}
\renewcommand{\P}{\mathbf{P}\xspace}
\newcommand{\NP}{\mathbf{NP}\xspace}
\newcommand{\PCP}{\mathbf{PCP}\xspace}
\newcommand{\F}{\mathbb{F}}

\begin{document}
\pagestyle{empty}

\begin{center}
%  {\small\scshape Università di Pisa -- Dipartimento di Matematica \par}
%  {\small Corso di Laurea Triennale in Matematica \par}
%  \vspace{0.8em}
  {\Large\bfseries Teorema PCP e inapprossimabilità \par}
\end{center}

\vspace{0.8em}

\noindent
\begin{minipage}[t]{0.47\textwidth}
  {\large Relatori: \par}
  \vspace{0.2em}
  {\bfseries\large Prof. Alessandro Panconesi \par}
  {\bfseries\large Prof.ssa Alessandra Caraceni \par}
\end{minipage}%
\hfill
\begin{minipage}[t]{0.47\textwidth}\raggedleft
  {\large Candidato: \par}
  \vspace{0.2em}
  {\bfseries\large Lorenzo Ferrari \par}
\end{minipage}

\vspace{0.8em}
\hrule height 0.4pt
\vspace{0.8em}

\begin{center}
  {\large\bfseries Abstract \par}
\end{center}
\vspace{0.2em}

Lo scopo di questa tesi è lo studio del Teorema PCP, risultato fondamentale
in informatica teorica, dimostrato in una prima versione nel 1992 da Arora e
Safra e completato lo stesso anno da Arora, Lund, Motwani, Sudan e Szegedy.

Una dimostrazione si legge tipicamente dall'inizio alla fine per assicurarsi
che tutti i passaggi logici siano validi, poiché un singolo errore potrebbe
invalidare l'intero argomento.
Immaginiamo ora ci vengano fornite mille pagine di dimostrazione di un teorema
molto complicato, per verificare la quale ci piacerebbe adottare la seguente
procedura molto sbrigativa: lette cento righe scelte a caso, accettiamo la
dimostrazione come valida se e soltanto se queste non contengono errori.
Un tale protocollo di verifica è impensabile per dimostrazioni in
formati classici nel caso di dimostrazioni fallaci, infatti, non c'è quasi nessuna
garanzia di correttezza.
Il Teorema PCP, tuttavia, afferma che esiste un modo di scrivere dimostrazioni
per cui questa verifica ha successo con alta probabilità.
Questa tesi intende presentare in maniera precisa il Teorema PCP, analizzarne le
conseguenze per l'inapprossimabilità di problemi di ottimizzazione combinatoria
e infine dimostrarlo.

Nel primo capitolo introduciamo le nozioni base della teoria della
complessità, a partire dalle classi $\P$ ed $\NP$, rispettivamente la classe
dei problemi che si risolvono efficientemente e la classe dei problemi le cui
soluzioni si verificano efficientemente.
Descriviamo i problemi $\NP$-completi $\mathtt{SAT}$, $\mathtt{3SAT}$,
$\mathtt{3COL}$ e soprattutto $\mathtt{QUADEQ}$: il problema di determinare se
un sistema di equazioni quadratiche su $\F_2$ ha soluzione.
Formalizziamo poi il modello di verifica probabilistica di dimostrazioni ed
enunciamo il Teorema PCP.
Come primo risultato significativo, dimostriamo una versione debole del
Teorema PCP, nota come Teorema PCP esponenziale.
La dimostrazione fa uso del codice di Walsh-Hadamard per rappresentare
dimostrazioni come funzioni lineari. Sarà utile anche il test di linearità di
Blum-Luby-Rubinfeld, analizzato mediante analisi di Fourier discreta
sull'ipercubo $(\mathbb{Z} / 2\mathbb{Z})^n$.

Nel secondo capitolo esploriamo il legame tra la verifica probabilistica
di dimostrazioni e l'inapprossimabilità di problemi di ottimizzazione
combinatoria. Generalizziamo $\mathtt{SAT}$, $\mathtt{3SAT}$ e $\mathtt{3COL}$
introducendo i \emph{Constraint Satisfaction Problems} ($\mathtt{CSP}$) e
mostriamo che il Teorema PCP ammette una formulazione equivalente
intrinsecamente legata all'inapprossimabilità di $\CSP$.
Come conseguenza diretta, non esiste un algoritmo efficiente in grado di
approssimare il valore di soddisfacibilità di istanze $\CSP$ oltre una certa
soglia costante a meno che $\P = \NP$.
Concludiamo con un risultato di inapprossimabilità per il problema della
cricca massima: non si può trovare efficientemente neanche una cricca grande
un miliardesimo della cricca massima, a meno che $\P = \NP$.

Il terzo capitolo è dedicato ai grafi expander e alle loro proprietà
pseudo-casuali. Per esempio, per ogni nodo $v$ di un expander con $n$ nodi,
la distribuzione del $k$-esimo nodo di una passeggiata aleatoria che parte da
$v$ è molto vicina a quella uniforme per $k = O(\log n)$.
Mostriamo inoltre che, per molte applicazioni, i nodi di una passeggiata
aleatoria di lunghezza $k$ funzionano bene quasi quanto $k$ nodi scelti a caso.
Ciò permette il cosiddetto ``riciclo di bit random'', che incluso nella
dimostrazione di inapprossimabilità per il problema della cricca massima
fornisce un risultato ancora migliore.

Nell'ultimo capitolo percorriamo la dimostrazione combinatoria
del Teorema PCP proposta da Irit Dinur nel 2006 basandoci sulla trattazione di
Arora e Barak.
La dimostrazione si basa sull'\emph{amplificazione del gap}, ottenuta grazie
all'impiego degli expander e alle loro proprietà pseudo-casuali.
Tra gli ingredienti per concludere la dimostrazione figura anche il Teorema PCP
esponenziale, dimostrato nel primo capitolo.

\end{document}
